\documentclass[11pt]{article}
\usepackage[polish]{babel}
\usepackage[T1]{fontenc}
\usepackage[utf8]{inputenc}
\usepackage{geometry}
\usepackage{graphicx}
\usepackage{booktabs}
\usepackage{amsmath}
\usepackage{tikz}
\usepackage{caption}
\usepackage{subcaption}
\usepackage{multirow}
\geometry{margin=2.5cm}

% --- biblioteki TikZ do rysunkow logiczno-czasowych (rysunki inline, bez plikow PNG) ---
\usetikzlibrary{shapes.gates.logic.US}
\usetikzlibrary{arrows.meta}
\usetikzlibrary{positioning}
\usetikzlibrary{calc}

\title{Projekt układu sumy liczby 2-bitowej z 3-bitową na bramkach NOR}
\author{}
\date{}

\begin{document}

\maketitle

\section{Podsumowanie i wnioski}

\subsection{Ocena zaprojektowanego układu}

Zaprojektowany układ sumy liczby 2-bitowej z 3-bitową sumą spełnia wszystkie stawiane przed nim wymagania funkcjonalne. Zastosowanie podejścia strukturalnego opartego na bramkach NOR okazało się wysoce efektywne. Układ działa w sposób deterministyczny i bezbłędnie wyznacza wynik sumy wraz z przeniesieniem najstarszego bitu.

W trakcie realizacji projektu wykonano:
\begin{itemize}
    \item analizę teoretyczną sumatora 2-bitowego z 3-bitową sumą,
    \item specyfikację wejść i wyjść,
    \item analizę logiczną funkcji wyjściowych,
    \item minimalizację funkcji na bramkach NOR,
    \item opracowanie schematów logicznych i tabelarycznych,
    \item weryfikację działania za pomocą symulacji i tablicy prawdy.
\end{itemize}

Układ charakteryzuje się prostą strukturą, którą można łatwo zminimalizować pod kątem liczby bramek oraz liczby połączeń.

\subsection{Korzyści wynikające z minimalizacji bramek NOR}

Minimalizacja funkcji logicznej przyniosła wymierne korzyści:
\begin{itemize}
    \item redukcja zasobów sprzętowych: liczba potrzebnych bramek spadła w stosunku do wariantu bezpośredniego,
    \item ograniczenie poboru prądu: mniejsza liczba przełączających się elementów wewnątrz struktur logicznych bezpośrednio przekłada się na mniejsze straty mocy,
    \item zwiększenie szybkości: skrócenie ścieżki krytycznej sygnału zminimalizowało opóźnienie propagacyjne, co pozwala na pracę układu przy wyższych częstotliwościach.
\end{itemize}

Realizacja wyłącznie na bramkach NOR pozwoliła na ujednolicenie struktury układu i uproszczenie jego interpretacji. W praktyce oznacza to mniejsze zapotrzebowanie na różne typy bramek logicznych oraz łatwiejszą realizację w strukturach cyfrowych.

Tabela porównawcza liczby bramek przed i po minimalizacji przedstawia się następująco:
\begin{table}[h]
\centering
\caption{PorÓwnanie liczby bramek przed i po minimalizacji}
\begin{tabular}{lll}
\toprule
Wariant & Liczba bramek & Uwagi \\
\midrule
Wariant początkowy (SOP) & większa liczba AND/OR/NOT & brak minimalizacji \\
Wariant zminimalizowany (NOR) & mniejsza liczba NOR & realizacja przyjazna dla NOR \\
\bottomrule
\end{tabular}
\end{table}

\subsection{Ograniczenia rozwiązania}

Prezentowany układ posiada ograniczenie architektoniczne -- realizuje wyłącznie funkcję sumy 2-bitowej z 3-bitową sumą. Nie dostarcza on informacji o relacji rzędu między liczbami, czyli nie wskazuje, która z liczb jest większa lub mniejsza. Dodatkowo, wraz ze wzrostem szerokości słowa danych (np. do 8 lub 16 bitów), struktura kaskadowa musiałaby zostać rozszerzona, co wymusiłoby zwiększenie liczby poziomów bramek i przyśpieszenie opóźnień.

W praktyce ograniczenia te są naturalnym skutkiem przyjętego zakresu projektu, który obejmował wyłącznie sumę dwóch liczb dwubitowych z wynikiem 3-bitowym.

\subsection{Możliwości rozwoju projektu}

Projekt może zostać rozbudowany w następujących kierunkach:
\begin{itemize}
    \item rozszerzenie funkcjonalności: dodanie bloków logicznych odpowiedzialnych za wykrywanie relacji większości i mniejszości,
    \item skalowalność: przekształcenie układu w sumator magistrali n-bitowej poprzez zastosowanie struktury kaskadowej lub drzewiastej,
    \item implementacja w języku opisu sprzętu: przeniesienie projektu na poziom opisu języka programowania układów programowalnych, co umożliwiłoby automatyczną syntezę w strukturach CPLD/FPGA.
\end{itemize}

Szczególnie interesującym kierunkiem rozwoju jest zastosowanie bramek NOR w strukturach programowalnych, gdzie możliwe jest wielokrotne użycie tych samych struktur bramkowych do realizacji różnych funkcji logicznych.

\subsection{Zakres pracy}

W pracy wykonano kompletną analizę projektu sumatora 2-bitowego z 3-bitową sumą na bramkach NOR. Praca obejmuje:
\begin{itemize}
    \item wstęp teoretyczny,
    \item opis techniczny projektu,
    \item minimalizację funkcji logicznej,
    \item schematy logiczne, tabelaryczne i wykresowe,
    \item symulację i testowanie,
    \item podsumowanie i wnioski.
\end{itemize}

Projekt został przygotowany w formie gotowej do dalszych prac rozwojowych oraz do wykorzystania w strukturach cyfrowych.

\begin{figure}[h]
\centering
\includegraphics[width=0.8\textwidth]{syg1.png}
\caption{Schemath końcowy układu}
\label{fig:schemat_koncowy}
\end{figure}

Na schemacie końcowym przedstawiono układ w formie uproszczonej, z wydzieleniem dwóch poziomów sumacyjnych oraz wyjścia przeniesienia. Rysunek \ref{fig:schemat_koncowy} przedstawia układ w formie hierarchicznej: wejścia są przetwarzane na poziomie młodszym, a następnie na poziomie starszym, a wynik końcowy jest sumą 3-bitową.

\begin{figure}[h]
\centering
\includegraphics[width=0.8\textwidth]{syg2.png}
\caption{Schemath logiczny układu}
\label{fig:schemat_log_konc}
\end{figure}

Rysunek \ref{fig:schemat_log_konc} przedstawia szczegółowo zależności między wejściami a wyjściami w formie bramek NOR.

\begin{figure}[h]
\centering
\includegraphics[width=0.8\textwidth]{syg3.png}
\caption{Końcowy rysunek podsumowujący}
\label{fig:koniec}
\end{figure}

Praca została zakończona. Wszystkie punkty zostały wykonane zgodnie z planem.

\end{document}
